% Edited conversion of source p. 31.
% The source line drawing is preserved as a base64-encoded asset and decoded by LuaLaTeX.
\directlua{
  dofile("scripts/base64_asset.lua")
  decode_base64_file(
    "figures/15-ribbed-shell/ribbed-shell.jpg.b64",
    "figures/15-ribbed-shell/ribbed-shell.generated.jpg"
  )
}

\ReportHeading
  {Про моделювання нестаціонарного деформування гнучкої тришарової циліндричної оболонки еліптичного перерізу з ребристим наповнювачем}
  {Є.~А. Сторожук, А.~В. Лисенко, І.~С. Чернишенко}
  {Інститут механіки ім. С.~П. Тимошенка НАН України}
  {}

Розглянуто тришарову циліндричну оболонку еліптичного перерізу з дискретним поздовжньо-поперечним ребристим наповнювачем за дії нестаціонарного навантаження. Неоднорідна тришарова пружна структура складається із зовнішньої та внутрішньої еліптичних циліндричних обшивок, жорстко з’єднаних поздовжніми й поперечними ребрами (рис.~1).

\begin{figure}[htbp]
 \centering
 \includegraphics[width=0.58\textwidth]{figures/15-ribbed-shell/ribbed-shell.generated.jpg}
 \caption{Тришарова циліндрична оболонка еліптичного перерізу з ребристим наповнювачем.}
\end{figure}

Розв’язувальні рівняння отримано із застосуванням геометрично нелінійної теорії оболонок і стрижнів типу Тимошенка у квадратичному наближенні~\cite{ribbed:golovko2012}. Зв’язок між внутрішніми зусиллями й моментами та деформаціями обшивок і ребер жорсткості визначено законом Гука для ортотропних матеріалів.

Рівняння коливань тришарової пружної структури з дискретним наповнювачем виведено з варіаційного принципу стаціонарності Гамільтона--Остроградського. Після підстановки виразів для потенціальної та кінетичної енергій обшивок і ребер отримано дві групи рівнянь --- для обшивок і для ребристого наповнювача.

Чисельний алгоритм ґрунтується на інтегро-інтерполяційному методі побудови різницевих схем за просторовими координатами та явній скінченно-різницевій схемі інтегрування за часом. Він включає: 1) визначення розв’язку у гладкій області обшивок; 2) визначення розв’язку на лінії розриву вздовж осі $i$-го ребра; 3) визначення розв’язку на лінії розриву вздовж осі $j$-го ребра.

Як числовий приклад досліджено вплив геометричної нелінійності та параметрів ребер на напружено-деформований стан тришарової циліндричної панелі еліптичного перерізу з поздовжнім ребристим наповнювачем. Поперечні краї панелі шарнірно закріплено, поздовжні залишено вільними; навантаження задано у вигляді поверхневого імпульсу.

\begin{thebibliography}{9}
\bibitem{ribbed:golovko2012}
К.~Г. Головко, П.~З. Луговой, В.~Ф. Мейш,
\emph{Динамика неоднородных оболочек при нестационарных нагрузках}, под ред. А.~Н. Гузя.
Киев: Издательско-полиграфический центр ``Киевский университет'', 2012, 541~с.
\end{thebibliography}
